% sec_intro.tex -- Section 1: Introduction
\section{Introduction}\label{sec:intro}

\subsection{The game}\label{ssec:intro-game}

Tic-tac-toe on the board $n^d=\{0,\dots,n-1\}^d$ is played by two players who alternately occupy an
empty cell. A \emph{line} is a set of $n$ cells of the form $\{(f_1(t),\dots,f_d(t)) : 0\le t\le n-1\}$,
where each coordinate $f_i(t)$ is constant, equal to $t$, or equal to $n-1-t$, and not all coordinates
are constant; there are $((n+2)^d-n^d)/2$ lines~\cite[Theorem~1]{GolombHales2002}. The player who first
occupies all $n$ cells of a line wins, and if the board is filled without a winner, the game is a draw.
The game is finite, so either one of the players has a winning strategy or both players can force a
draw; by the strategy-stealing argument the second player never has a winning strategy
(\cref{prop:stealing}). For each pair $(n,d)$ the question is therefore whether the first player wins
or the game is a draw.

A draw is usually proved by proving more. In the \emph{Maker--Breaker game} on the same board the
players are called Maker and Breaker; Maker wins by occupying all cells of some line, and Breaker wins
otherwise, that is, by occupying at least one cell of every line. If Breaker wins this game as the
second player, then the same strategy prevents the first player of the ordinary game from winning, and
the ordinary game is a draw; in the terminology of Beck~\cite{Beck2008} it is then even a \emph{Strong
Draw}. See~\cite{Beck2008,HefetzKrivelevichStojakovicSzabo2014,Krivelevich2014} for background on
positional games. The simplest strategies for Breaker are \emph{pairing strategies}: Breaker fixes two
cells on every line such that the pairs of different lines are disjoint, and whenever Maker takes a
cell of a pair, Breaker takes the other one. A pairing needs two cells for every line, so it can exist
only if the number of cells is at least twice the number of lines.

\subsection{Known results}\label{ssec:intro-known}

Hales and Jewett~\cite{HalesJewett1963} proved that the $n^d$ game is a draw by a pairing strategy if
$n\ge 3^d-1$ and $n$ is odd, or $n\ge 2^{d+1}-2$ and $n$ is even, and that for every fixed $n$ the first
player wins if $d$ is large enough; see~\cite[p.~300]{ErdosSelfridge1973}
and~\cite[Theorem~2.1 and Corollary~3.4]{Krivelevich2014}. For $d=3$ these bounds give a draw only for
even $n\ge 14$ and odd $n\ge 27$, and much more is known (\cref{tab:known}).

\begin{table}[ht]
\centering\small
\begin{tabular}{@{}cccl@{}}
\toprule
$n$ & lines & cells & outcome of the $n\times n\times n$ game\\
\midrule
$\le 3$ & & & first player wins (elementary)\\
$4$ & $76$ & $64$ & first player wins \cite{Patashnik1980,AllisSchoo1992,Allis1994}\\
$5$ & $109$ & $125$ & unknown; conjectured to be a draw \cite{Beck2008,GolombHales2002}\\
$6$ & $148$ & $216$ & draw (\cref{thm:main})\\
$7$ & $193$ & $343$ & draw (\cite{CainRaymondKallersjo2025}, preprint; \cref{cor:seven-intro})\\
$\ge 8$ & $3n^2+6n+4$ & $n^3$ & draw by a pairing strategy \cite{GolombHales2002}\\
\bottomrule
\end{tabular}
\caption{Tic-tac-toe on the board $n\times n\times n$.}
\label{tab:known}
\end{table}

The first player wins for $n\le 4$: this is elementary for $n\le 3$, and for $n=4$, the game Qubic, it
is a computer-assisted theorem of Patashnik~\cite{Patashnik1980}, which was later reproved with
different search methods by Allis and Schoo~\cite{AllisSchoo1992,Allis1994}. For $n\ge 8$ the second
player draws by an explicit pairing strategy, which goes back to
Golomb~\cite[p.~181]{GolombHales2002}; see also~\cite[Chapter~22]{BerlekampConwayGuy2003}. For $n\le 7$
there is no pairing strategy on the
empty board, because the $3n^2+6n+4$ lines would need more cells than there are: $218>125$ for $n=5$,
$296>216$ for $n=6$, and $386>343$ for $n=7$. The cases $n=5,6,7$ have long been expected to be draws.
Golomb and Hales~\cite[Table~1]{GolombHales2002} place them, marked as uncertain, in the region where
the second player can force a draw but not by a pairing strategy, and Beck~\cite[p.~1]{Beck2008} calls
$5\times5\times5$ tic-tac-toe ``a hopeless open problem (it is conjectured to be a draw game)''. The
criterion of Erd\H{o}s and Selfridge~\cite{ErdosSelfridge1973} does not decide any of the three cases
(\cref{ssec:remarks}).

Recently Cain, Raymond and K\"allersj\"o~\cite{CainRaymondKallersjo2025} gave, in a preprint, a
computer-assisted proof that Breaker wins the Maker--Breaker game on the $7\times7\times7$ board: after
three or four prescribed moves of Breaker a pairing strategy exists against every continuation of
Maker. They list $n=5$ and $n=6$ as the two remaining open cases. For $n=6$ they sketch a
plan~\cite[Section~6.2]{CainRaymondKallersjo2025}: with the first four moves Breaker occupies a cell on
each of the four space diagonals; then ``[i]t is possible that Breaker could, within the next few
rounds of play, eliminate enough additional winning lines to establish a viable pairing strategy, while
making sure to block any potential threats from Maker'', and ``[w]ith sufficient care beyond the scope
of what we are willing to attempt, it may be possible to resolve the $6^3$ game computationally''. The
present paper carries out this plan.

\emph{Computer-assisted results on other positional games.} Many games have been solved with the help
of computers; see the survey~\cite{vandenHerikUiterwijkvanRijswijck2002}. Patashnik's solution of Qubic
combined a search with expert knowledge supplied by hand~\cite{Patashnik1980}; see~\cite[pp.~95--96]{Allis1994}.
Go-Moku was solved by Allis, van den Herik and Huntjens~\cite{AllisvandenHerikHuntjens1996} with
threat-space search and proof-number search~\cite{AllisvanderMeulenvandenHerik1994}. The threats, the
forced replies and the AND--OR search of the present paper belong to the same circle of ideas; here,
however, the search only serves to find a certificate, which is then checked independently. This
division of labor between an untrusted search and small programs that check its output is that of
certifying algorithms~\cite{McConnellEtAl2011} and of large proofs by satisfiability solving in which
solving and verifying are separate steps~\cite{HeuleKullmannMarek2016}. Closer to
our leaf conditions are proofs of draws that replace search by an argument of matching type:
Uiterwijk's ``set matching''~\cite{Uiterwijk2019} is an enhancement of the Hales--Jewett pairing
strategy for positions of $k$-in-a-row games. (For a recent computer-assisted solution of games of
this family see~\cite{HsuEtAl2020}, where the $(7,7,5)$-game and the $(8,8,5)$-game are treated.) What
is specific to the $6\times6\times6$ board is described in \cref{ssec:intro-method,ssec:intro-seven}.

\subsection{The result}\label{ssec:intro-result}

From now on the board is $\Q=\{0,\dots,5\}^3$, with $216$ cells and $148$ lines, and Maker is the
player who moves first. A \emph{round} consists of a move of Maker followed by a move of Breaker.

% the main theorem must not be broken across pages
\needlines{7}
\begin{theorem}[computer-assisted]\label{thm:main}\interlinepenalty=10000
In the Maker--Breaker game on the $148$ lines of the $6\times6\times6$ board, \Breaker, moving second,
has a winning strategy: a strategy that prescribes an empty cell whenever \Breaker{} is to move and under
which \Maker{} never owns all six cells of a line. Consequently, ordinary $6\times6\times6$ tic-tac-toe
is a draw: each of the two players has a strategy with which this player never loses.
\end{theorem}

The theorem is computer-assisted in the following precise sense. \Cref{sec:board,sec:pairing,sec:leaf,sec:game}
define a finite \emph{abstract game} and the notion of a \emph{valid certificate} for it
(\cref{def:certificate}): a table of moves, a list of leaves and a bound $\Rmax$ on the number of rounds
that are played from the table. They prove, in the text and without a computer: \emph{if a valid certificate exists, then both
assertions of \cref{thm:main} hold} (\cref{thm:soundness,thm:sigma} and \cref{cor:draw}; the formal proof
of \cref{thm:main} is given in \cref{thm:main-proof}). That the two certificate files distributed with
this paper form a valid certificate with $\Rmax=10$ is \cref{fact:H} in \cref{sec:certificate}. It is a
finite statement about two explicit files, which was established by running programs, and it is the only
computed statement on which \cref{thm:main} depends.

By a standard argument, a player who wins a Maker--Breaker game when moving second also wins it when
moving first~\cite{Krivelevich2014}; so \Breaker{} wins the game of \cref{thm:main} also as the first
player. We do not use this.

To our knowledge, \cref{thm:main} is the first proof that the $6\times6\times6$ game is a draw; we
found no earlier proof or claim in the literature. A consequence of \cref{thm:main} is a new proof of
the result of the preprint~\cite{CainRaymondKallersjo2025} for the $7\times7\times7$ board.

\begin{corollary}[computer-assisted, via \cref{thm:main}]\label{cor:seven-intro}
In the Maker--Breaker game on the $193$ lines of the $7\times7\times7$ board, \Breaker, moving second,
has a winning strategy. Consequently, ordinary $7\times7\times7$ tic-tac-toe is a draw.
\end{corollary}

This is \cref{cor:seven} in \cref{sec:discussion}. Its proof uses the decomposition of the
$7\times7\times7$ board by its three mid-planes, after Golomb and Hales~\cite{GolombHales2002}: the
cells outside the mid-planes form a $6\times6\times6$ board, on which Breaker plays the strategy of
\cref{thm:main}, and the lines inside the mid-planes admit a pairing, which is given explicitly and can
be checked without a computer. The proof does not depend on the computation of the
preprint~\cite{CainRaymondKallersjo2025}.
Granting \cref{thm:main} and the result for $n=7$, the $5\times5\times5$ game is, to our knowledge, the
only $n\times n\times n$ game whose outcome is unknown.

\subsection{The method in outline}\label{ssec:intro-method}

\emph{No pairing on the empty board.} The $148$ lines would need $296$ cells, and there are $216$. A
stone of Breaker kills the lines through its cell, at most seven of them, and a count shows that no
pairing of the surviving lines can exist before Breaker has placed seven stones (\cref{prop:seven}).
Seven or eight well-chosen stones do admit a pairing on an otherwise empty board
(\cref{ex:T7,ex:T8} and \cref{fact:T-pairings}). When Breaker has placed them, however, Maker has placed
as many stones, the pairing has to avoid these, and whether this is possible depends on where they are.

\emph{Admissible positions and leaf conditions.} Call a line \emph{live} if it contains no stone of
Breaker. Two configurations leave Breaker no choice: a live line with five stones of Maker (a
\emph{threat}), and two live lines, each with four stones of Maker, that meet in an empty cell (a
\emph{danger pair}). Breaker answers them at once, by
taking the sixth cell of the threat or the common cell of the danger pair. A set of stones of Maker
without such configurations is called \emph{admissible} (\cref{def:admissible}); admissibility is exactly
Hall's condition for one and for two lines. A \emph{leaf condition} is a condition on Breaker's
stones $B$, on a known subset $F$ of Maker's stones and on the number $t$ of Maker's stones which
guarantees that a pairing exists for \emph{every} admissible placement of Maker's remaining $t-|F|$
stones. We use a sufficient condition, $\LEAF(B,F,t)$, which is derived from Hall's theorem and is
decided by exact computations with augmenting paths (\cref{sec:leaf}). It requires a pairing that
avoids the known stones $F$ and, in addition, that every connected set of at least three live lines has
at least $t-|F|$ cells more than a pairing needs. The reason why this suffices is short: a set of lines
that violates Hall's condition can be taken connected; one or two lines cannot violate it when Maker's
stones are admissible; and a connected set of at least three lines with that many spare cells cannot be
spoiled by the $t-|F|$ stones whose positions are not known (\cref{thm:leaf-implies-C}).

\emph{Forced replies and lazy revelation.} Before a leaf condition holds, Breaker plays from a table.
A table with an entry for every position would be far too large. The table used here is indexed by the
triple $(B,F,t)$ alone, where $F$ is the set of \emph{revealed} stones of Maker. A stone of Maker is
revealed in two situations only: when it takes part in a threat or in a danger pair, which Breaker has
to answer by the forced reply, and when the table tells Breaker to occupy a cell that turns out to be
occupied, in which case the stone found there is revealed and the table is consulted again. A move of
Maker that creates neither a threat nor a danger pair thus leads to a single successor, wherever it was
made. This is made precise by the abstract game of \cref{sec:game}; \cref{fig:round} shows one round
of it, and \cref{tab:opening} lists the table entries for the first three rounds. A \emph{certificate}
consists of a
table of moves and a list of leaves; it is \emph{valid} if every state that can be reached is a leaf
that satisfies the leaf condition or has all the entries it needs (\cref{def:certificate}). The
soundness theorem
(\cref{thm:soundness}) shows that a valid certificate yields a strategy for Breaker under which Maker's
stones are admissible at the beginning of every round, until a leaf is reached and a pairing strategy
takes over. One more point has to be taken care of: a single move of Maker might create a threat and a
danger pair, or two danger pairs, which cannot be answered by one reply. This requires at least eleven
stones of Maker (\cref{prop:lost-stones}) and therefore cannot happen during the first ten rounds.

\emph{The certificate.} The certificate of this paper lets Breaker play from the table for at most ten
rounds. It has $93{,}164$ table entries and $106{,}723$ leaves, up to the $1{,}152$ symmetries of the
board, and Breaker switches to a pairing strategy after round $8$, $9$ or $10$ (\cref{sec:certificate}).
It was found by a heuristic search in about one minute. The search is not part of the proof: the
certificate was verified from scratch by programs that were written independently of the search and of
each other (\cref{sec:verification}; the sense of ``independently'' is explained in
\cref{ssec:impl} and in the Declarations at the end of the paper).

\subsection{Comparison with the proof for the \texorpdfstring{$7\times7\times7$}{7x7x7} board}
\label{ssec:intro-seven}

The idea of letting Breaker place a few stones and then switch to a pairing found by bipartite matching
is that of the preprint~\cite{CainRaymondKallersjo2025}, and for the $6\times6\times6$ board the plan of opening on
the space diagonals and blocking threats on the way to a pairing is sketched
there (\cref{ssec:intro-known}). On the $7\times7\times7$ board three or four opening moves
suffice. Maker has then placed so few stones that every line still has at least three empty cells, so
no threat has to be answered before the pairing takes over, and the pairings are computed for every
position that can have arisen~\cite{CainRaymondKallersjo2025}. On the $6\times6\times6$ board the
situation is different in three respects. At least seven stones of Breaker are needed before any
pairing can exist, so Maker has time to build threats and danger pairs, which must be answered on the
way; the states of our certificate have $143{,}498$ successors that arise from a threat and $28{,}275$
that arise from a danger pair (\cref{tab:tree}). The number of positions that Maker can reach in eight
to ten rounds is far too large for an enumeration, which is why the table depends on Maker's stones only through the
revealed ones and why the pairing has to be guaranteed by a leaf condition that covers all admissible placements
of the others. Finally, the pairing is reached only after eight to ten rounds, along a strategy tree whose
validity is the computed part of the proof.

\subsection{The nature of the proof}\label{ssec:intro-nature}

We state plainly what is proved in the text, without a computer, what is computed, and what has to be
trusted.

\emph{Written proofs.} All statements called Theorem, Proposition, Lemma or Corollary are proved in the
paper. Three of them are marked ``computer-assisted'': \cref{thm:main} and the two statements of its
corollary for the $7\times7\times7$ board (\cref{cor:seven-intro,cor:seven}) depend, in addition to
the written proofs, on \cref{fact:H}. The written proofs include
the geometry of the board (\cref{lem:geometry,lem:group}), the pairing strategy, the implication from
the leaf condition to the existence of pairings (\cref{thm:leaf-implies-C}), the soundness of the
abstract game (\cref{thm:soundness,thm:sigma}) and the strategy-stealing argument. The only results
quoted from the literature in these proofs are Hall's theorem, in the form of the max-flow min-cut
theorem with integrality, and, for the second assertion of \cref{thm:main}, the determinacy of finite
games with perfect information.

\emph{Computed.} Mathematical statements that we establish by running programs and state formally are
called \emph{Computed Facts} and name the programs that established them.
\Cref{sec:certificate,sec:verification,sec:repro} and the appendices, those of the supplementary material
included, also report computational results in running text and in tables, each with its source; none
of these is used in the proof. Exactly one
Computed Fact, \cref{fact:H}, is used in the proof of
\cref{thm:main}: it says that the two certificate files form a valid certificate. \Cref{sec:verification}
states the two algorithms that check this, precisely enough to be implemented from the paper, and
proves them correct (\cref{prop:verify-correct,prop:leaf-correct}); the main verifier uses a variant of
the second algorithm, which is justified in \cref{lem:deleted}. The check was carried out by the two
verifiers of the supplementary archive, which were written separately, and by two further programs
written separately during a review of the project (\cref{ssec:impl}), all with the same result and the
same counts of states in every round. During an internal refereeing of the first version of this paper
three more programs were written from the text of the paper: two complete verifiers, which accept the
certificate, and a transcription of the first algorithm without the leaf test; all three reproduce the
same counts (\cref{ssec:impl,ssec:own}). All these programs use integer
arithmetic only; no floating-point computation and no random choice enters the proof. The other
Computed Facts and the cross-checks reported in \cref{sec:verification}, among them an integer program
solved in floating-point arithmetic and one million simulated games, are not used in the proof.

\emph{Trusted.} A reader who accepts \cref{thm:main} trusts the written proofs and at least one
implementation of the verification algorithms, together with the compiler or interpreter, the
operating system and the hardware on which it ran (\cref{ssec:trust}). Put differently, \cref{thm:main}
can fail only if a written proof in \cref{sec:board,sec:pairing,sec:leaf,sec:game} is wrong, or if the
algorithms of \cref{sec:verification} do not check \cref{def:certificate}, or if every one of the
implementations has one of the defects listed in \cref{ssec:what}. The programs, including all
verifiers, were written with AI-based tools; see the Declarations at the end of the paper and
\cref{ssec:impl} for the sense in which they are independent. The search that produced the
certificate is not trusted. Nothing has been formally verified in a proof assistant; we come back to
this in \cref{ssec:open}.

\subsection{Availability}\label{ssec:intro-artifact}

The certificate, the sources of all programs and the outputs of the runs reported here are collected
in a supplementary archive, which is described in \cref{sec:repro}; see ``Data and code
availability'' in the Declarations at the end of the paper. The certificate consists of two text files,
whose formats are specified in
\cref{sec:formats}, with the SHA-256 checksums
\begin{center}\small
\begin{tabular}{@{}ll@{}}
\fn{strategy.tbl} & \texttt{13629393fcdcee87a03f2c9c73f640d9584d2134a46f90dad9393dc92c9ce3fc}\\
\fn{leaves.txt}   & \texttt{3c14c7814e296e812b4fdf3cda2f446c6ce4865d15d351980e7fecc7690fd470}
\end{tabular}
\end{center}
Verifying the certificate with the supplied programs takes a few seconds with the main verifier and a
few minutes with the second one (\cref{ssec:certifying}).

\subsection{Organization of the paper}\label{ssec:intro-org}

The paper consists of a main part and of supplementary material, both in this file. The main part
contains everything that the proof of \cref{thm:main} uses and everything that is needed to replicate
the verification of the certificate. \Cref{sec:board} describes the board, its $148$ lines and its group of $1{,}152$
symmetries. \Cref{sec:pairing} treats pairings and Hall's condition, and \cref{sec:leaf} the admissible
sets and the leaf conditions. \Cref{sec:game} defines the abstract game, certificates and Breaker's
strategy, and proves the soundness theorem and, from it, \cref{thm:main}. \Cref{sec:certificate} states
the computed fact and describes the certificate. \Cref{sec:verification} gives the verification
algorithms with their correctness proofs and reports on the verifiers, and \cref{sec:repro} says how to
verify the certificate with the programs of the archive or with a program of one's own.
\Cref{sec:discussion} contains the corollary on the $7\times7\times7$ board, remarks and open problems;
it is followed by the Declarations. \Cref{sec:formats} specifies the file formats and the canonical form
and gives test vectors for an independent verifier, and \cref{sec:seven-pairing} gives the pairing used
for the $7\times7\times7$ board. The notation is collected in \cref{tab:notation}.

The supplementary material follows the references and begins on page~\pageref{part:supp}. Its
appendices are numbered S1 to S7. They contain results that are not used in the proof of
\cref{thm:main} (\cref{sec:extra,sec:extra-game}), the description of the search that found the
certificate (\cref{sec:supp-search}), the step-by-step description of the leaf test of the main verifier
and the details of the audits and of the cross-checks (\cref{sec:supp-checks}), the integer program that
was used as a cross-check (\cref{sec:ip}), the details of the archive and of the reproduction of the
computation (\cref{sec:supp-repro}), and further tables (\cref{sec:supp-tables}). None of them is needed
for the proof or for a replication. The supplementary material is part of this paper; it is not to be
confused with the supplementary archive of programs and data (\cref{ssec:intro-artifact,sec:repro}).
